\documentclass[10pt,aspectratio=169]{beamer}
\usepackage{../../Common/course_slides}
\usepackage{../../Common/course_macros}
\usepackage[spanish,provide=*]{babel}

\title[Semana 1: Vectores, Rectas y Planos]{Geometría del Espacio Euclidiano: Vectores, Rectas y Planos}
\subtitle{De la Geometría de Flechas a la Estructura de Espacio Vectorial en $\R^3$}
\author{Geometría y Física I: Mecánica Clásica y Formas de Cartan}
\institute{}
\date{}

\begin{document}
\setlength{\abovedisplayskip}{3pt}
\setlength{\belowdisplayskip}{3pt}

% ---------------------------------------------------------------------
% SLIDE 1: PORTADA
% ---------------------------------------------------------------------
\begin{frame}
  \titlepage
\end{frame}

% ---------------------------------------------------------------------
% SLIDE 2: PREGUNTAS DE ACTIVACIÓN GEOMÉTRICA
% ---------------------------------------------------------------------
\begin{frame}{La Naturaleza Geométrica del Espacio Físico}
  \vspace{0.1cm}
  Al describir el movimiento, el equilibrio o las interacciones en el mundo físico:
  \vspace{0.3cm}
  \begin{enumerate}\itemsep14pt
    \item Si las coordenadas $(x, y, z)$ de un objeto dependen de dónde elija el observador colocar su origen, ¿por qué magnitudes como la velocidad, la aceleración y las fuerzas no cambian al trasladar el sistema de referencia?
    \item ¿Qué diferencia estructural existe entre una simple colección o lista de números y un objeto geométrico dotado de magnitud, dirección y reglas propias de transformación?
    \item ¿Por qué la suma de dos posiciones en el espacio carece de significado geométrico, mientras que la resta de dos posiciones define un vector de desplazamiento genuino?
  \end{enumerate}
\end{frame}

% ---------------------------------------------------------------------
% SLIDE 3: DEDUCCIÓN GUIADA - DESPLAZAMIENTO
% ---------------------------------------------------------------------
\begin{frame}{Invarianza del Desplazamiento y Acción de Vectores}
  \begin{block}{Desplazamiento Relativo entre Puntos}
    \footnotesize
    Una partícula parte de la posición $P_0 = (1, 2, 1)\,\text{m}$ y después de $1\,\text{s}$ pasa por $P_1 = (3, 5, 3)\,\text{m}$ moviéndose en línea recta a velocidad constante.
    \begin{enumerate}
      \itemsep2pt
      \item[(a)] \textbf{La suma de posiciones:} ¿Tiene sentido físico o geométrico «sumar el lugar $P_0$ con el lugar $P_1$»? ¿Qué le pasa a $P_0 + P_1$ si un segundo observador mueve el origen a $(10, 0, 0)$?
      \item[(b)] \textbf{El vector desplazamiento:} Calcula $\vec{v} = P_1 - P_0$. Demuestra que si trasladamos el origen sumando una constante $\vec{c} \in \R^3$ a todos los puntos ($P' = P + \vec{c}$), el vector $\vec{v}$ \textbf{no cambia}.
      \item[(c)] \textbf{Construcción de la recta:} Si la partícula mantiene su dirección $\vec{v}$ en todo instante $t \in \R$, escribe una fórmula para su posición $\vec{r}(t)$ y determina en qué instante $t$ alcanza la altura $z = 9\,\text{m}$.
    \end{enumerate}
  \end{block}
\end{frame}

% ---------------------------------------------------------------------
% SLIDE 4: DEFINICIÓN FORMAL - ESPACIO AFÍN DE WEYL CON TIKZ
% ---------------------------------------------------------------------
\begin{frame}{El Espacio Afín: Axiomática de Weyl}
  \begin{columns}[c]
    \begin{column}{0.54\textwidth}
      \begin{block}{Definición de Espacio Afín $(A, V, +)$}
        \scriptsize
        Un conjunto de puntos $A$ y un espacio vectorial $V$ con una traslación $+: A \times V \to A$ que satisface:
        \begin{enumerate}
          \itemsep1pt
          \item $P + \vec{0} = P$ para todo $P \in A$.
          \item $(P + \vec{u}) + \vec{v} = P + (\vec{u} + \vec{v})$.
          \item \textbf{Vector entre dos puntos:} Para cada par $(P, Q)$, existe un \textbf{único} vector $Q - P \in V$ tal que $P + (Q - P) = Q$.
        \end{enumerate}
      \end{block}
      \vspace{0.05cm}
      \scriptsize
      \textbf{Relación de Chasles:} $(Q - P) + (R - Q) = R - P$.\\[2pt]
      \textbf{Principio Físico:} No hay origen universal; el espacio es homogéneo.
    \end{column}
    \begin{column}{0.44\textwidth}
      \centering
      \begin{tikzpicture}[scale=0.82]
        \coordinate (P) at (0.3, 0.4);
        \coordinate (Q) at (3.2, 0.9);
        \coordinate (R) at (1.8, 2.8);
        \coordinate (O) at (-0.7, -0.6);

        % Chasles triangle
        \draw[-{Stealth}, line width=1.5pt, color=orangeacc] (P) -- (Q) node[midway, below=2pt, font=\scriptsize\bfseries] {$\vec{u} = Q - P$};
        \draw[-{Stealth}, line width=1.5pt, color=navyaccent] (Q) -- (R) node[midway, right=2pt, font=\scriptsize\bfseries] {$\vec{v} = R - Q$};
        \draw[-{Stealth}, line width=1.7pt, color=blackmain] (P) -- (R) node[midway, above left=1pt, font=\scriptsize\bfseries] {$R - P = \vec{u}+\vec{v}$};

        % Points
        \fill[blackmain] (P) circle (2.2pt) node[below left=1pt, font=\scriptsize\bfseries] {$P$};
        \fill[blackmain] (Q) circle (2.2pt) node[below right=1pt, font=\scriptsize\bfseries] {$Q$};
        \fill[blackmain] (R) circle (2.2pt) node[above=2pt, font=\scriptsize\bfseries] {$R$};

        % Dashed origin
        \draw[dashed, color=graymuted, -{Stealth}] (O) -- (P) node[pos=0.5, left=1pt, font=\tiny] {$\vec{r}_P$};
        \draw[dashed, color=graymuted, -{Stealth}] (O) -- (Q) node[pos=0.5, right=1pt, font=\tiny] {$\vec{r}_Q$};
        \fill[graymuted] (O) circle (1.8pt) node[below left=1pt, font=\tiny\bfseries] {$O$};
      \end{tikzpicture}
      \vspace{0.05cm}

      {\scriptsize Identidad de Chasles e independencia del origen $O$.}
    \end{column}
  \end{columns}
\end{frame}

% ---------------------------------------------------------------------
% SLIDE 5: DEFINICIÓN FORMAL - ESPACIO VECTORIAL
% ---------------------------------------------------------------------
\begin{frame}{El Espacio Vectorial Real: Definición Formal y Axiomas}
  \begin{columns}[c]
    \begin{column}{0.54\textwidth}
      \footnotesize
      \begin{itemize}
        \item \textbf{Espacio Afín vs. Espacio Vectorial:} Un punto $P$ es una ubicación; la diferencia $\vec{v} = P_1 - P_0$ es un \textbf{vector} en un espacio con neutro $\vec{0}$.
        \item \textbf{Definición Formal $(V, +, \cdot)$ sobre $\R$:} Conjunto cerrado bajo suma $\vec{u}+\vec{v}$ y producto escalar $a\vec{u}$ ($a \in \R$) con \textbf{8 axiomas}:
        \begin{enumerate}
          \scriptsize
          \item Conmutatividad: $\vec{u} + \vec{v} = \vec{v} + \vec{u}$.
          \item Asociatividad: $(\vec{u}+\vec{v})+\vec{w} = \vec{u}+(\vec{v}+\vec{w})$.
          \item Neutro aditivo único $\vec{0}$: $\vec{v} + \vec{0} = \vec{v}$.
          \item Inverso aditivo $-\vec{v}$: $\vec{v} + (-\vec{v}) = \vec{0}$.
          \item Asociatividad escalar: $a(b\vec{v}) = (ab)\vec{v}$ y $1\vec{v} = \vec{v}$.
          \item Distributividades: $a(\vec{u}+\vec{v}) = a\vec{u}+a\vec{v}$, \ $(a+b)\vec{v} = a\vec{v}+b\vec{v}$.
        \end{enumerate}
      \end{itemize}
    \end{column}
    \begin{column}{0.43\textwidth}
      \centering
      \begin{tikzpicture}[scale=0.82]
        \coordinate (O) at (0,0);
        \coordinate (U) at (3.4, 0.3);
        \coordinate (V) at (1.1, 2.3);
        \coordinate (S) at ($(U)+(V)$);
        \fill[orangeacc!8] (O) -- (U) -- (S) -- (V) -- cycle;
        \draw[dashed, color=graymuted] (U) -- (S) (V) -- (S);
        \draw[-{Stealth}, line width=1.4pt, color=blackmain] (O) -- (U) node[midway, below=2pt, font=\scriptsize\bfseries] {$\vec{u}$};
        \draw[-{Stealth}, line width=1.4pt, color=blackmain] (O) -- (V) node[midway, left=2pt, font=\scriptsize\bfseries] {$\vec{v}$};
        \draw[-{Stealth}, line width=1.6pt, color=orangeacc] (O) -- (S) node[above right=0pt, font=\scriptsize\bfseries] {$\vec{u}+\vec{v}$};
        \draw[-{Stealth}, line width=1.4pt, color=navyaccent] (V) -- (U) node[pos=0.68, above right=2pt, font=\scriptsize\bfseries] {$\vec{u}-\vec{v}$};
      \end{tikzpicture}
      \vspace{0.1cm}

      {\scriptsize Las dos diagonales: $\vec{u}+\vec{v}$ (suma) y $\vec{u}-\vec{v}$ (resta).}
    \end{column}
  \end{columns}
\end{frame}

% ---------------------------------------------------------------------
% SLIDE 6: PREGUNTAS DEDICADAS - ESPACIO AFÍN Y VECTORES
% ---------------------------------------------------------------------
\begin{frame}{Espacio Afín y la Naturaleza del Vector}
  \vspace{0.2cm}
  \small
  \begin{enumerate}\itemsep14pt
    \item En un espacio afín, cualquier punto $P$ puede conectarse con otro $Q$ mediante un único vector $\vec{v} = Q - P$. ¿Cómo formaliza esta relación la noción de que el espacio físico es homogéneo y no posee un centro absoluto?
    \item Si fijamos un origen arbitrario $O$, cada punto $P$ se identifica con un vector posición $\vec{r}_P = P - O$. ¿Qué ocurre con esta correspondencia biunívoca si cambiamos el origen a $O' = O + \vec{d}$?
    \item ¿Por qué magnitudes vectoriales físicas como el momento lineal $\vec{p} = m\dot{\vec{r}}$ y el campo eléctrico $\vec{E}$ se suman siguiendo la regla del paralelogramo, pero sus puntos de aplicación pertenecen a un espacio afín?
  \end{enumerate}
\end{frame}

% ---------------------------------------------------------------------
% SLIDE 7: DE LA LEY DE COSENOS AL PRODUCTO PUNTO CON TIKZ
% ---------------------------------------------------------------------
\begin{frame}{De la Ley de Cosenos al Producto Punto}
  \begin{columns}[c]
    \begin{column}{0.54\textwidth}
      \footnotesize
      \textbf{¿Cómo medir el ángulo $\theta$ entre dos vectores?}\\[4pt]
      En el triángulo formado por $\vec{u}$, $\vec{v}$ y la resta $\vec{u}-\vec{v}$, la \textbf{Ley de los Cosenos} establece:
      \[
      \|\vec{u} - \vec{v}\|^2 = \|\vec{u}\|^2 + \|\vec{v}\|^2 - 2\|\vec{u}\|\|\vec{v}\|\cos\theta
      \]
      Desarrollando en componentes cartesianas en $\R^3$:
      \[
      \sum_{i=1}^3 (u_i - v_i)^2 = \sum_{i=1}^3 (u_i^2 + v_i^2) - 2\sum_{i=1}^3 u_i v_i
      \]
      Los términos cuadráticos se cancelan y emerge el \textbf{producto punto}:
      \[
      \boxed{\vec{u}\cdot\vec{v} \equiv \sum_{i=1}^3 u_i v_i = \|\vec{u}\|\|\vec{v}\|\cos\theta}
      \]
    \end{column}
    \begin{column}{0.44\textwidth}
      \centering
      \begin{tikzpicture}[scale=0.88]
        \coordinate (O) at (0,0);
        \coordinate (U) at (3.2, 0.4);
        \coordinate (V) at (1.5, 2.3);

        \draw[-{Stealth}, line width=1.5pt, color=blackmain] (O) -- (U) node[midway, below=2pt, font=\scriptsize\bfseries] {$\vec{u}$};
        \draw[-{Stealth}, line width=1.5pt, color=blackmain] (O) -- (V) node[midway, left=2pt, font=\scriptsize\bfseries] {$\vec{v}$};
        \draw[-{Stealth}, line width=1.6pt, color=orangeacc] (V) -- (U) node[midway, above right=1pt, font=\scriptsize\bfseries] {$\vec{w} = \vec{u} - \vec{v}$};

        \pic[draw=navyaccent, fill=navyaccent!15, angle radius=0.65cm, "$\theta$", font=\scriptsize\bfseries] {angle = U--O--V};

        \fill[blackmain] (O) circle (2pt) node[below left=1pt, font=\scriptsize] {$O$};
      \end{tikzpicture}
      \vspace{0.1cm}

      {\scriptsize El triángulo fundamental: $\|\vec{u}-\vec{v}\|^2$ conecta el álgebra de coordenadas con el ángulo $\theta$.}
    \end{column}
  \end{columns}
\end{frame}

% ---------------------------------------------------------------------
% SLIDE 8: AXIOMAS DEL PRODUCTO INTERNO Y MÉTRICAS CON TIKZ
% ---------------------------------------------------------------------
\begin{frame}{Axiomas del Producto Interno y Métricas Ponderadas}
  \begin{columns}[c]
    \begin{column}{0.54\textwidth}
      \begin{block}{Definición de Producto Interno $\langle \vec{u}, \vec{v} \rangle$}
        \scriptsize
        Una operación bilineal $\langle \cdot, \cdot \rangle: V \times V \to \R$ que satisface:
        \begin{enumerate}
          \itemsep1pt
          \item \textbf{Simetría:} $\langle \vec{u}, \vec{v} \rangle = \langle \vec{v}, \vec{u} \rangle$.
          \item \textbf{Bilinealidad:} $\langle c\vec{u} + \vec{w}, \vec{v} \rangle = c\langle \vec{u}, \vec{v} \rangle + \langle \vec{w}, \vec{v} \rangle$.
          \item \textbf{Positividad:} $\langle \vec{u}, \vec{u} \rangle \ge 0$, con igualdad $\iff \vec{u} = \vec{0}$.
        \end{enumerate}
      \end{block}
      \vspace{0.05cm}
      \scriptsize
      \textbf{Métricas Ponderadas (Medios Anisótropos):}\\[2pt]
      En física, las direcciones pueden tener distinta rigidez:
      \[
      \langle \vec{u}, \vec{v} \rangle_{\text{aniso}} = u_1 v_1 + 4 u_2 v_2
      \]
      La «circunferencia unitaria» se deforma en una elipse.
    \end{column}
    \begin{column}{0.44\textwidth}
      \centering
      \begin{tikzpicture}[scale=0.88]
        % Canonical circle
        \draw[color=graymuted, thin, -{Stealth}] (-1.6,0) -- (1.6,0) node[right, font=\tiny] {$x_1$};
        \draw[color=graymuted, thin, -{Stealth}] (0,-1.6) -- (0,1.6) node[above, font=\tiny] {$x_2$};
        \draw[line width=1.3pt, color=navyaccent] (0,0) circle (1.1cm);
        \node[color=navyaccent, font=\scriptsize\bfseries] at (0, -1.35) {$\|\vec{x}\|_{\text{can}} = 1$};

        % Anisotropic ellipse
        \begin{scope}[xshift=2.9cm]
          \draw[color=graymuted, thin, -{Stealth}] (-1.6,0) -- (1.6,0) node[right, font=\tiny] {$x_1$};
          \draw[color=graymuted, thin, -{Stealth}] (0,-1.6) -- (0,1.6) node[above, font=\tiny] {$x_2$};
          \draw[line width=1.5pt, color=orangeacc] (0,0) ellipse (1.1cm and 0.55cm);
          \node[color=orangeacc, font=\scriptsize\bfseries] at (0, -1.35) {$x_1^2 + 4x_2^2 = 1$};
        \end{scope}
      \end{tikzpicture}
      \vspace{0.05cm}

      {\scriptsize Geometría isotrópica (círculo) vs. medio anisótropo (elipse métrica).}
    \end{column}
  \end{columns}
\end{frame}

% ---------------------------------------------------------------------
% SLIDE 9: ESPACIOS NORMADOS Y DEFINICIÓN DE ÁNGULO
% ---------------------------------------------------------------------
\begin{frame}{Espacios Normados y la Definición Geométrica de Ángulo}
  \begin{block}{Definición Formal de Norma}
    \footnotesize
    Una \textbf{norma} en $V$ es una función $\|\cdot\|: V \to \R_{\ge 0}$ que asigna a cada vector su longitud satisfaciendo:
    \begin{enumerate}
      \itemsep1pt
      \item \textbf{No-degeneración:} $\|\vec{u}\| = 0 \iff \vec{u} = \vec{0}$.
      \item \textbf{Homogeneidad absoluta:} $\|c\vec{u}\| = |c|\|\vec{u}\|$ para todo $c \in \R$ y $\vec{u} \in V$.
      \item \textbf{Desigualdad triangular:} $\|\vec{u} + \vec{v}\| \le \|\vec{u}\| + \|\vec{v}\|$.
    \end{enumerate}
  \end{block}
  \vspace{0.08cm}
  \begin{columns}[c]
    \begin{column}{0.51\textwidth}
      \footnotesize
      \textbf{Norma Inducida por el Producto Interno:}
      \[
      \boxed{\|\vec{u}\| \equiv \sqrt{\langle \vec{u}, \vec{u} \rangle}}
      \]
      Todo producto interno define intrínsecamente la distancia entre puntos: $d(P, Q) = \|Q - P\|$.
    \end{column}
    \begin{column}{0.46\textwidth}
      \footnotesize
      \textbf{Definición Intrínseca del Ángulo:}
      \[
      \boxed{\cos\theta \equiv \frac{\langle \vec{u}, \vec{v} \rangle}{\|\vec{u}\|\|\vec{v}\|}, \quad \theta \in [0, \pi]}
      \]
      \textbf{Pregunta crucial:} ¿Cómo garantizamos que este cociente siempre viva en $[-1, 1]$?
    \end{column}
  \end{columns}
\end{frame}

% ---------------------------------------------------------------------
% SLIDE 10: TEOREMA DE CAUCHY-SCHWARZ CON TIKZ
% ---------------------------------------------------------------------
\begin{frame}{Teorema Fundamental: La Desigualdad de Cauchy-Schwarz}
  \begin{columns}[c]
    \begin{column}{0.55\textwidth}
      \begin{block}{Enunciado del Teorema}
        \footnotesize
        Para cualesquiera vectores $\vec{u}, \vec{v}$ en un espacio con producto interno:
        \[
        \boxed{|\langle \vec{u}, \vec{v} \rangle| \le \|\vec{u}\|\|\vec{v}\|}
        \]
        La igualdad se alcanza si y solo si $\vec{u}$ y $\vec{v}$ son linealmente dependientes ($\vec{u} = c\vec{v}$).
      \end{block}
      \vspace{0.1cm}
      \footnotesize
      \textbf{Garantía Geométrica:}
      \[
      -1 \le \frac{\langle \vec{u}, \vec{v} \rangle}{\|\vec{u}\|\|\vec{v}\|} \le 1
      \]
      Asegura que el ángulo $\theta$ entre cualquier par de vectores está siempre matemáticamente bien definido.
    \end{column}
    \begin{column}{0.43\textwidth}
      \centering
      \begin{tikzpicture}[scale=0.88]
        \coordinate (O) at (0,0);
        \coordinate (V) at (3.6,0);
        \coordinate (U) at (2.4, 1.8);
        \coordinate (P) at (2.4,0);

        \draw[-{Stealth}, line width=1.3pt, color=graymuted] (O) -- (V) node[right, font=\scriptsize\bfseries] {$\vec{v}$};
        \draw[-{Stealth}, line width=1.5pt, color=blackmain] (O) -- (U) node[midway, above left=1pt, font=\scriptsize\bfseries] {$\vec{u}$};
        \draw[-{Stealth}, line width=2pt, color=orangeacc] (O) -- (P) node[midway, below=3pt, font=\scriptsize\bfseries] {$\frac{|\langle\vec{u},\vec{v}\rangle|}{\|\vec{v}\|}$};
        \draw[dashed, color=navyaccent, thick] (P) -- (U);

        \draw[color=navyaccent, thin] (2.15,0) -- (2.15,0.25) -- (2.4,0.25);
        \pic[draw=orangeacc, angle radius=0.6cm, "$\theta$", font=\scriptsize] {angle = V--O--U};
      \end{tikzpicture}
      \vspace{0.08cm}

      {\scriptsize La proyección ortogonal (cateto) nunca puede superar a la hipotenusa $\|\vec{u}\|$.}
    \end{column}
  \end{columns}
\end{frame}

% ---------------------------------------------------------------------
% SLIDE 11: DEMOSTRACIÓN DE CAUCHY-SCHWARZ (SLIDE DEDICADA)
% ---------------------------------------------------------------------
\begin{frame}{Demostración de la Desigualdad de Cauchy-Schwarz}
  \footnotesize
  Sea $\vec{v} \neq \vec{0}$ (si $\vec{v} = \vec{0}$, la desigualdad $0 \le 0$ es trivial). Consideremos la combinación $\vec{u} + \lambda\vec{v}$ para cualquier escalar $\lambda \in \R$.
  \vspace{0.15cm}

  Por la propiedad de positividad definida del producto interno, la longitud al cuadrado nunca es negativa:
  \[
  P(\lambda) \equiv \|\vec{u} + \lambda \vec{v}\|^2 = \langle \vec{u} + \lambda \vec{v},\; \vec{u} + \lambda \vec{v} \rangle \ge 0 \qquad (\forall \lambda \in \R)
  \]
  Expandiendo por bilinealidad y simetría obtenemos un polinomio cuadrático en $\lambda$:
  \[
  P(\lambda) = \|\vec{u}\|^2 + 2\lambda\langle \vec{u}, \vec{v} \rangle + \lambda^2 \|\vec{v}\|^2 \ge 0
  \]
  Como el coeficiente principal es $A = \|\vec{v}\|^2 > 0$ y la parábola $P(\lambda)$ nunca cruza por debajo del eje horizontal ($P(\lambda) \ge 0$), la ecuación $P(\lambda) = 0$ tiene a lo más una raíz real.
  \vspace{0.15cm}

  Por tanto, su discriminante $\Delta = B^2 - 4AC$ no puede ser estrictamente positivo:
  \[
  \Delta = (2\langle \vec{u}, \vec{v} \rangle)^2 - 4\|\vec{v}\|^2 \|\vec{u}\|^2 \le 0 \implies 4\langle \vec{u}, \vec{v} \rangle^2 \le 4\|\vec{u}\|^2 \|\vec{v}\|^2
  \]
  Dividiendo entre $4$ y extrayendo la raíz cuadrada positiva a ambos miembros:
  \[
  \boxed{|\langle \vec{u}, \vec{v} \rangle| \le \|\vec{u}\|\|\vec{v}\|} \qquad \blacksquare
  \]
\end{frame}

% ---------------------------------------------------------------------
% SLIDE 12: TEOREMA DE LA DESIGUALDAD TRIANGULAR CON TIKZ
% ---------------------------------------------------------------------
\begin{frame}{Teorema Fundamental: La Desigualdad Triangular}
  \begin{columns}[c]
    \begin{column}{0.54\textwidth}
      \begin{block}{Enunciado del Teorema}
        \footnotesize
        Para cualesquiera vectores $\vec{u}, \vec{v}$ en un espacio con producto interno y norma inducida:
        \[
        \boxed{\|\vec{u} + \vec{v}\| \le \|\vec{u}\| + \|\vec{v}\|}
        \]
        La igualdad se alcanza si y solo si $\vec{u} = c\vec{v}$ con $c \ge 0$ (vectores colineales en el mismo sentido).
      \end{block}
      \vspace{0.1cm}
      \footnotesize
      \textbf{Significado Físico y Geométrico:}\\[2pt]
      En cualquier triángulo, la distancia en línea recta entre el inicio y el final es estrictamente menor o igual que la suma de los tramos individuales (el rodeo).
    \end{column}
    \begin{column}{0.44\textwidth}
      \centering
      \begin{tikzpicture}[scale=0.92]
        \coordinate (O) at (0,0);
        \coordinate (U) at (2.8, 0.4);
        \coordinate (V) at (1.2, 2.0);
        \coordinate (S) at ($(U)+(V)$);

        \draw[-{Stealth}, line width=1.5pt, color=navyaccent] (O) -- (U) node[midway, below=2pt, font=\scriptsize\bfseries] {$\vec{u}$};
        \draw[-{Stealth}, line width=1.5pt, color=navyaccent] (U) -- (S) node[midway, right=2pt, font=\scriptsize\bfseries] {$\vec{v}$};
        \draw[-{Stealth}, line width=2pt, color=orangeacc] (O) -- (S) node[midway, above left=1pt, font=\scriptsize\bfseries] {$\vec{u}+\vec{v}$};

        \fill[blackmain] (O) circle (2pt) node[below left=1pt, font=\tiny] {Inicio};
        \fill[blackmain] (S) circle (2pt) node[above right=1pt, font=\tiny] {Fin};
      \end{tikzpicture}
      \vspace{0.08cm}

      {\scriptsize El camino directo $\|\vec{u}+\vec{v}\|$ frente al camino poligonal $\|\vec{u}\|+\|\vec{v}\|$.}
    \end{column}
  \end{columns}
\end{frame}

% ---------------------------------------------------------------------
% SLIDE 13: DEMOSTRACIÓN DE LA DESIGUALDAD TRIANGULAR (SLIDE DEDICADA)
% ---------------------------------------------------------------------
\begin{frame}{Demostración de la Desigualdad Triangular}
  \footnotesize
  Calculamos el cuadrado de la norma del vector suma $\vec{u}+\vec{v}$ expandiendo por bilinealidad y simetría:
  \[
  \|\vec{u} + \vec{v}\|^2 = \langle \vec{u} + \vec{v},\; \vec{u} + \vec{v} \rangle = \|\vec{u}\|^2 + 2\langle \vec{u}, \vec{v} \rangle + \|\vec{v}\|^2
  \]
  Dado que para cualquier número real $x \le |x|$, acotamos el término cruzado por su valor absoluto:
  \[
  \|\vec{u} + \vec{v}\|^2 \le \|\vec{u}\|^2 + 2|\langle \vec{u}, \vec{v} \rangle| + \|\vec{v}\|^2
  \]
  Aplicando directamente el \textbf{Teorema de Cauchy-Schwarz} ($|\langle \vec{u}, \vec{v} \rangle| \le \|\vec{u}\|\|\vec{v}\|$):
  \[
  \|\vec{u} + \vec{v}\|^2 \le \|\vec{u}\|^2 + 2\|\vec{u}\|\|\vec{v}\| + \|\vec{v}\|^2
  \]
  El miembro derecho es exactamente el desarrollo del trinomio cuadrado perfecto:
  \[
  \|\vec{u} + \vec{v}\|^2 \le (\|\vec{u}\| + \|\vec{v}\|)^2
  \]
  Puesto que tanto $\|\vec{u}+\vec{v}\|$ como $\|\vec{u}\|+\|\vec{v}\|$ son cantidades no negativas, extraemos la raíz cuadrada a ambos lados conservando el orden:
  \[
  \boxed{\|\vec{u} + \vec{v}\| \le \|\vec{u}\| + \|\vec{v}\|} \qquad \blacksquare
  \]
\end{frame}

% ---------------------------------------------------------------------
% SLIDE 14: DEDUCCIÓN GUIADA - PROYECCIÓN ORTOGONAL
% ---------------------------------------------------------------------
\begin{frame}{Descomposición de Fuerzas y Proyección Ortogonal}
  \begin{block}{Descomposición Ortogonal sobre una Dirección Fija}
    \footnotesize
    Considera una fuerza $\vec{F}$ que empuja un objeto guiado por un riel recto con vector director $\vec{v} \neq \vec{0}$. Queremos descomponer $\vec{F}$ en dos partes mutuamente perpendiculares:
    \[
    \vec{F} = \vec{F}_\parallel + \vec{F}_\perp, \qquad \vec{F}_\parallel = c\vec{v}, \quad \vec{F}_\perp \cdot \vec{v} = 0.
    \]
    \begin{enumerate}
      \itemsep2pt
      \item[(a)] Toma el producto escalar de la ecuación con $\vec{v}$: $\vec{F}\cdot\vec{v} = c(\vec{v}\cdot\vec{v}) + \vec{F}_\perp\cdot\vec{v}$. Usando la ortogonalidad de $\vec{F}_\perp$, despeja el escalar $c$.
      \item[(b)] Escribe la fórmula vectorial explícita de $\vec{F}_\parallel = \text{proy}_{\vec{v}}(\vec{F})$ y de la fuerza normal al riel $\vec{F}_\perp$.
      \item[(c)] Verifica que $\|\vec{F}\|^2 = \|\vec{F}_\parallel\|^2 + \|\vec{F}_\perp\|^2$ calculando directamente $\langle \vec{F}_\parallel + \vec{F}_\perp, \vec{F}_\parallel + \vec{F}_\perp \rangle$.
    \end{enumerate}
  \end{block}
\end{frame}

% ---------------------------------------------------------------------
% SLIDE 15: TEORÍA Y TIKZ - PROYECCIONES
% ---------------------------------------------------------------------
\begin{frame}{El Teorema de Descomposición Ortogonal y Proyecciones}
  \begin{columns}[c]
    \begin{column}{0.53\textwidth}
      \footnotesize
      \begin{itemize}
        \item \textbf{Proyección Ortogonal sobre $\vec{v}$:}
        \[
        \boxed{\text{proy}_{\vec{v}}(\vec{u}) = \left(\frac{\vec{u}\cdot\vec{v}}{\|\vec{v}\|^2}\right)\vec{v}}
        \]
        \item \textbf{Componente Perpendicular (Normal):}
        \[
        \boxed{\vec{u}_\perp = \vec{u} - \text{proy}_{\vec{v}}(\vec{u}), \qquad \vec{u}_\perp \cdot \vec{v} = 0}
        \]
        \item \textbf{Ortogonalidad Pitagórica:}
        \[
        \|\vec{u}\|^2 = \|\text{proy}_{\vec{v}}(\vec{u})\|^2 + \|\vec{u}_\perp\|^2
        \]
      \end{itemize}
    \end{column}
    \begin{column}{0.44\textwidth}
      \centering
      \begin{tikzpicture}[scale=0.82]
        \foreach \x in {0.4, 1.0, 1.6, 2.2, 2.8} {
          \draw[-{Stealth}, color=orangeacc!60, thin, densely dashed] (\x, 3.1) -- (\x, 2.6);
        }
        \node[color=orangeacc, font=\scriptsize\itshape] at (1.6, 3.35) {Rayos perpendiculares al riel};

        \draw[color=graymuted, line width=1pt] (-0.3,0) -- (4.8,0);

        \coordinate (O) at (0,0);
        \coordinate (U) at (3.0, 2.2);
        \coordinate (P) at (3.0, 0);
        \coordinate (V) at (4.5, 0);

        \draw[color=orangeacc, thick] (2.75, 0) -- (2.75, 0.25) -- (3.0, 0.25);
        \draw[dashed, color=orangeacc, line width=1.1pt] (P) -- (U) node[midway, right=2pt, font=\scriptsize\bfseries] {$\vec{u}_\perp = \vec{u} - \vec{u}_\parallel$};

        \draw[-{Stealth}, line width=1.3pt, color=blackmain] (O) -- (V) node[above=1pt, font=\scriptsize\bfseries] {$\vec{v}$};
        \draw[-{Stealth}, line width=2.2pt, color=orangeacc] (O) -- (P) node[midway, below=3pt, font=\scriptsize\bfseries] {$\vec{u}_\parallel = \text{proy}_{\vec{v}}(\vec{u})$};
        \draw[-{Stealth}, line width=1.6pt, color=navyaccent] (O) -- (U) node[midway, above left=1pt, font=\scriptsize\bfseries] {$\vec{u}$};

        \pic[draw=navyaccent, fill=navyaccent!10, angle radius=0.6cm, "$\theta$", font=\scriptsize] {angle = P--O--U};
      \end{tikzpicture}
    \end{column}
  \end{columns}
\end{frame}

% ---------------------------------------------------------------------
% SLIDE 16: PREGUNTAS DEDICADAS - GEOMETRÍA MÉTRICA
% ---------------------------------------------------------------------
\begin{frame}{Geometría Métrica y Proyecciones}
  \vspace{0.2cm}
  \small
  \begin{enumerate}\itemsep14pt
    \item ¿Por qué la proyección ortogonal $\vec{u}_\parallel = \text{proy}_{\vec{v}}(\vec{u})$ es el vector en la dirección de $\vec{v}$ que minimiza la distancia $\|\vec{u} - \lambda\vec{v}\|$ para todo $\lambda \in \R$?
    \item ¿Cómo nos permite el signo algebraico de $\vec{u}\cdot\vec{v}$ clasificar el ángulo relativo (agudo si $>0$, perpendicular si $=0$, obtuso si $<0$) sin evaluar funciones trigonométricas?
    \item ¿Por qué la descomposición $\vec{u} = \vec{u}_\parallel + \vec{u}_\perp$ es única y no depende de la orientación de los ejes cartesianos auxiliares?
  \end{enumerate}
\end{frame}

% ---------------------------------------------------------------------
% SLIDE 17: TEORÍA Y TIKZ - RECTA
% ---------------------------------------------------------------------
\begin{frame}{La Ecuación Vectorial de la Recta}
  \begin{columns}[c]
    \begin{column}{0.48\textwidth}
      \footnotesize
      \begin{itemize}
        \item \textbf{Ecuación Vectorial de la Recta:}
        \[
        \boxed{\vec{r}(t) = \vec{p}_0 + t\vec{v}, \qquad t \in \R, \quad \vec{v} \neq \vec{0}}
        \]
        \item \textbf{Subespacio Vectorial 1D vs. Recta Afín:}
        \begin{itemize}
          \scriptsize
          \item Si $\vec{p}_0 = \vec{0}$, el conjunto $W = \{t\vec{v} \mid t \in \R\}$ contiene a $\vec{0}$ y es un \textbf{subespacio vectorial de dimensión 1}.
          \item Si $\vec{p}_0 \neq \vec{0}$, la recta $\mathcal{L} = \vec{p}_0 + W$ es una \textbf{variedad lineal afín} trasladada.
        \end{itemize}
        \item \textbf{Vector Tangente a una Curva:} La derivada $\dot{\vec{r}}(t) = \vec{v}$ representa la velocidad instantánea del movimiento rectilíneo.
      \end{itemize}
    \end{column}
    \begin{column}{0.49\textwidth}
      \centering
      \begin{tikzpicture}[scale=0.82]
        \draw[-{Stealth}, color=graymuted] (-0.3,0) -- (5.8,0) node[right, font=\scriptsize] {$x$};
        \draw[-{Stealth}, color=graymuted] (0,-0.3) -- (0,3.8) node[above, font=\scriptsize] {$y$};

        \coordinate (O) at (0,0);
        \coordinate (P0) at (1.6, 1.3);
        \coordinate (V) at (1.5, 0.7);
        \coordinate (P1) at ($(P0)+(V)$);
        \coordinate (Pt) at ($(P0)+2.1*(V)$);

        \draw[dashed, color=navyaccent!50, line width=1pt] ($(P0)-0.7*(V)$) -- ($(P0)+2.5*(V)$);
        \draw[color=navyaccent, line width=1.5pt] ($(P0)-0.3*(V)$) -- ($(P0)+2.2*(V)$);

        \draw[-{Stealth}, thick, color=blackmain] (O) -- (P0) node[pos=0.55, below right=-1pt, font=\scriptsize\bfseries] {$\vec{p}_0$};
        \draw[-{Stealth}, thick, color=navyaccent] (O) -- (Pt) node[pos=0.62, below right=2pt, font=\scriptsize\bfseries] {$\vec{r}(t) = \vec{p}_0 + t\vec{v}$};

        \draw[-{Stealth}, line width=1.6pt, color=orangeacc] (P0) -- (P1) node[midway, above left=1pt, font=\scriptsize\bfseries] {$\vec{v}$};
        \draw[-{Stealth}, dashed, color=orangeacc, line width=1.2pt] (P1) -- (Pt) node[midway, above left=1pt, font=\scriptsize\bfseries] {$t\vec{v}$};

        \fill[blackmain] (P0) circle (2.2pt) node[above left=2pt, font=\scriptsize\bfseries] {$P_0$};
        \fill[orangeacc] (P1) circle (2pt) node[below right=1pt, font=\scriptsize] {$t=1$};
        \fill[navyaccent] (Pt) circle (2.2pt) node[above right=1pt, font=\scriptsize\bfseries] {$P(t)$};
      \end{tikzpicture}
    \end{column}
  \end{columns}
\end{frame}

% ---------------------------------------------------------------------
% SLIDE 18: DEDUCCIÓN GUIADA - PRODUCTO CRUZ
% ---------------------------------------------------------------------
\begin{frame}{Ortogonalidad Simultánea y Producto Vectorial}
  \begin{block}{Construcción de una Dirección Perpendicular Conjunta}
    \footnotesize
    Dos bordes coplanares de una placa triangular están dados por los vectores $\vec{u} = (u_1, u_2, u_3)$ y $\vec{v} = (v_1, v_2, v_3)$. Buscamos un vector no nulo $\vec{n} = (n_1, n_2, n_3)$ que sea simultáneamente ortogonal a ambos bordes:
    \[
    \vec{u}\cdot\vec{n} = u_1 n_1 + u_2 n_2 + u_3 n_3 = 0, \qquad \vec{v}\cdot\vec{n} = v_1 n_1 + v_2 n_2 + v_3 n_3 = 0.
    \]
    \begin{enumerate}
      \itemsep2pt
      \item[(a)] Multiplica la primera ecuación por $v_3$ y la segunda por $u_3$, y réstalas para eliminar $n_3$. Obtén la relación entre $n_1$ y $n_2$: $(u_1 v_3 - u_3 v_1)n_1 + (u_2 v_3 - u_3 v_2)n_2 = 0$.
      \item[(b)] Repite para las otras parejas y deduce la solución canónica $\vec{n} = (u_2 v_3 - u_3 v_2,\; u_3 v_1 - u_1 v_3,\; u_1 v_2 - u_2 v_1)$.
      \item[(c)] ¿Qué sucede si intercambias el orden de los vectores ($\vec{v}$ por $\vec{u}$)? ¿Por qué la orientación de la normal con signo es indispensable al calcular el flujo a través de una superficie?
    \end{enumerate}
  \end{block}
\end{frame}

% ---------------------------------------------------------------------
% SLIDE 19: TEORÍA Y TIKZ - PRODUCTO CRUZ
% ---------------------------------------------------------------------
\begin{frame}{El Producto Cruz: Magnitud, Área y Orientación}
  \begin{columns}[c]
    \begin{column}{0.53\textwidth}
      \footnotesize
      \begin{itemize}
        \item \textbf{Definición Algebraica:}
        \[
        \vec{u} \times \vec{v} = (u_2 v_3 - u_3 v_2,\; u_3 v_1 - u_1 v_3,\; u_1 v_2 - u_2 v_1)
        \]
        \item \textbf{Magnitud = Área del Paralelogramo:}
        \[
        \|\vec{u}\times\vec{v}\| = \|\vec{u}\|\|\vec{v}\|\sin\theta = \text{Área}(\vec{u},\vec{v})
        \]
        \item \textbf{Antisimetría (Orientación / Mano Derecha):}
        \[
        \boxed{\vec{u}\times\vec{v} = -\vec{v}\times\vec{u} \implies \vec{u}\times\vec{u} = \vec{0}}
        \]
        \item \textbf{Propiedad Clave:} El signo codifica la orientación relativa del par ordenado $(\vec{u}, \vec{v})$.
      \end{itemize}
    \end{column}
    \begin{column}{0.44\textwidth}
      \centering
      \begin{tikzpicture}[scale=0.76]
        \coordinate (O) at (0,0);
        \coordinate (U) at (3.4, -0.25);
        \coordinate (V) at (1.4, 1.45);
        \coordinate (S) at ($(U)+(V)$);
        \coordinate (N) at (0, 2.8);
        \coordinate (NM) at (0, -1.4);

        \fill[orangeacc!12] (O) -- (U) -- (S) -- (V) -- cycle;
        \draw[dashed, color=orangeacc, thick] (U) -- (S) (V) -- (S);
        \node[color=blackmain, font=\scriptsize\bfseries] at (2.75, 0.62) {Área $= \|\vec{u}\times\vec{v}\|$};

        \draw[-{Stealth}, line width=1.5pt, color=navyaccent] (O) -- (U) node[midway, below=2pt, font=\scriptsize\bfseries] {$\vec{u}$};
        \draw[-{Stealth}, line width=1.5pt, color=blackmain] (O) -- (V) node[above left=1pt, font=\scriptsize\bfseries] {$\vec{v}$};
        \draw[-{Stealth}, line width=1.8pt, color=orangeacc] (O) -- (N) node[above, font=\scriptsize\bfseries] {$\vec{n} = \vec{u}\times\vec{v}$};
        \draw[-{Stealth}, line width=1.2pt, dashed, color=graymuted] (O) -- (NM) node[below, font=\scriptsize] {$\vec{v}\times\vec{u} = -\vec{u}\times\vec{v}$};

        \draw[-{Stealth}, color=orangeacc, thick] (0.95, -0.07) arc (-4:46:0.95);
        \node[color=orangeacc, font=\scriptsize\bfseries] at (1.15, 0.32) {$\theta$};
      \end{tikzpicture}
    \end{column}
  \end{columns}
\end{frame}

% ---------------------------------------------------------------------
% SLIDE 20: PREGUNTAS DEDICADAS - ORIENTACIÓN Y PRODUCTO CRUZ
% ---------------------------------------------------------------------
\begin{frame}{Orientación Espacial y Álgebra del Producto Cruz}
  \vspace{0.2cm}
  \small
  \begin{enumerate}\itemsep14pt
    \item ¿Por qué el producto vectorial $\vec{u}\times\vec{v}$ genera un vector perpendicular al plano generado por $\vec{u}$ y $\vec{v}$ únicamente en $\R^3$, a diferencia del producto escalar que existe en cualquier $\R^n$?
    \item ¿Por qué la propiedad de antisimetría $\vec{u}\times\vec{v} = -\vec{v}\times\vec{u}$ es la manifestación algebraica indispensable de la regla de la mano derecha y la orientación espacial?
    \item Si la magnitud $\|\vec{u}\times\vec{v}\|$ mide el área del paralelogramo, ¿por qué el momento de una fuerza (torque) $\vec{\tau} = \vec{r}\times\vec{F}$ y el momento angular $\vec{L} = \vec{r}\times\vec{p}$ se construyen naturalmente con este producto?
  \end{enumerate}
\end{frame}

% ---------------------------------------------------------------------
% SLIDE 21: DEDUCCIÓN GUIADA - PLANO
% ---------------------------------------------------------------------
\begin{frame}{Geometría del Plano: Descripción Paramétrica e Implícita}
  \begin{block}{Parametrización Bidimensional vs. Ecuación Normal}
    \footnotesize
    Queremos describir una superficie plana $\Pi$ que pasa por $P_0 = (1,0,2)$ y es generada por los vectores $\vec{v}_1 = (1, 1, 0)$ y $\vec{v}_2 = (0, 1, 1)$.
    \begin{enumerate}
      \itemsep2pt
      \item[(a)] \textbf{Descripción Interna (Paramétrica e Independencia Lineal):} Usamos dos coordenadas $(u,v) \in \R^2$ para movernos sobre el plano: $\vec{r}(u,v) = \vec{p}_0 + u\vec{v}_1 + v\vec{v}_2$. ¿Qué ocurriría geométricamente si $\vec{v}_2$ fuera múltiplo escalar de $\vec{v}_1$?
      \item[(b)] \textbf{Descripción Externa (Ecuación Normal):} Podemos caracterizar el mismo plano mediante una sola ecuación que restrinja la dirección perpendicular. Calcula $\vec{n} = \vec{v}_1 \times \vec{v}_2$ y escribe $\vec{n}\cdot(\vec{r} - \vec{p}_0) = 0$ en la forma $n_1 x + n_2 y + n_3 z = d$.
      \item[(c)] \textbf{Familia de Planos Paralelos:} ¿Qué representa geométricamente variar $d$ en $x - y + z = d$ para $d = 0, 1, 2, 3, \dots$? Si una partícula experimenta el desplazamiento $\vec{w} = (2, -1, 3)$, ¿cuántos planos de esta familia atraviesa?
    \end{enumerate}
  \end{block}
\end{frame}

% ---------------------------------------------------------------------
% SLIDE 22: TEORÍA Y TIKZ - PLANOS Y FOLIACIONES
% ---------------------------------------------------------------------
\begin{frame}{Planos en $\R^3$: Subespacios de Dimensión 2 y Foliaciones}
  \begin{columns}[c]
    \begin{column}{0.51\textwidth}
      \footnotesize
      \begin{itemize}
        \item \textbf{Independencia Lineal, Base y Dimensión:}
        $\{\vec{v}_1, \vec{v}_2\}$ son \emph{linealmente independientes} ($c_1\vec{v}_1+c_2\vec{v}_2=\vec{0} \implies c_1=c_2=0$) y forman una \textbf{base} del subespacio $W = \text{span}\{\vec{v}_1,\vec{v}_2\}$ ($\dim W = 2$).
        \item \textbf{Dualidad del Plano $\Pi = \vec{p}_0 + W$:}
        \begin{itemize}
          \scriptsize
          \item \textbf{Paramétrica ($2$ parámetros):} $\vec{r}(u,v) = \vec{p}_0 + u\vec{v}_1 + v\vec{v}_2$.
          \item \textbf{Normal ($1$ restricción):} $\vec{n}\cdot(\vec{r}-\vec{p}_0) = 0 \iff n_1 x + n_2 y + n_3 z = c$.
        \end{itemize}
        \item \textbf{Foliación por Planos Paralelos:}
        Al variar $c \in \R$, la ecuación $\vec{n}\cdot\vec{r} = c$ forma una familia de planos paralelos. El producto $\vec{n}\cdot\vec{w}$ mide cuántas láminas perfora $\vec{w}$.
      \end{itemize}
    \end{column}
    \begin{column}{0.46\textwidth}
      \centering
      \begin{tikzpicture}[scale=0.65]
        \foreach \y/\lbl/\col in {0/{c=1}/navyaccent!10, 1.05/{c=2}/navyaccent!18, 2.1/{c=3}/navyaccent!26, 3.15/{c=4}/navyaccent!34} {
          \fill[\col] (0,\y) -- (3.6,\y) -- (4.5,\y+0.75) -- (0.9,\y+0.75) -- cycle;
          \draw[color=navyaccent, thin] (0,\y) -- (3.6,\y) -- (4.5,\y+0.75) -- (0.9,\y+0.75) -- cycle;
          \node[right=2pt, font=\scriptsize, color=navyaccent] at (4.5, \y+0.5) {$\vec{n}\cdot\vec{r} = \lbl$};
        }

        \coordinate (A) at (1.5, 0.35);
        \coordinate (B) at (2.7, 3.55);
        \draw[-{Stealth}, line width=2pt, color=orangeacc] (A) -- (B);
        \node[above=2pt, font=\scriptsize\bfseries, color=orangeacc] at (2.4, 3.95) {$\Delta c = \vec{n}\cdot\vec{w} = 3$};

        \foreach \t in {0, 0.333, 0.666, 1.0} {
          \fill[blackmain] ($(A)!\t!(B)$) circle (2.2pt);
        }

        \draw[-{Stealth}, line width=1.4pt, color=blackmain] (0.65, 0.35) -- (0.65, 1.65) node[left=1pt, font=\scriptsize\bfseries] {$\vec{n}$};
      \end{tikzpicture}
    \end{column}
  \end{columns}
\end{frame}

% ---------------------------------------------------------------------
% SLIDE 23: DEDUCCIÓN GUIADA - INTERSECCIÓN
% ---------------------------------------------------------------------
\begin{frame}{Intersección de Planos y Sistemas de Restricciones}
  \begin{block}{Determinación de la Recta de Intersección}
    \footnotesize
    Dos planos $\Pi_1$ y $\Pi_2$ tienen vectores normales $\vec{n}_1 = (1, 0, -1)$ y $\vec{n}_2 = (0, 1, 1)$, con ecuaciones $x - z = 1$ y $y + z = 2$.
    \begin{enumerate}
      \item[(a)] \textbf{Conteo de Dimensión:} Cada plano impone $1$ restricción lineal independiente en $\R^3$. ¿Qué dimensión tiene la intersección $\Pi_1 \cap \Pi_2$?
      \item[(b)] \textbf{Dirección de la Recta:} Si el vector director $\vec{v}$ de la recta $\mathcal{L} = \Pi_1 \cap \Pi_2$ vive dentro de $\Pi_1$ y dentro de $\Pi_2$, ¿cómo es $\vec{v}$ respecto a $\vec{n}_1$ y $\vec{n}_2$? Usa el producto cruz para encontrar $\vec{v}$ de inmediato.
    \end{enumerate}
  \end{block}
\end{frame}

% ---------------------------------------------------------------------
% SLIDE 24: TEORÍA Y TIKZ - INTERSECCIÓN Y DISTANCIA
% ---------------------------------------------------------------------
\begin{frame}{Geometría de la Intersección y Distancia a un Plano}
  \begin{columns}[c]
    \begin{column}{0.56\textwidth}
      \footnotesize
      \begin{itemize}
        \item \textbf{Distancia de un Punto $Q$ a un Plano $\Pi$:}
        Proyección de $\vec{q}-\vec{p}_0$ sobre la normal unitaria $\hat{n}$:
        \[
        \boxed{D(Q,\Pi) = \frac{|\vec{n}\cdot(\vec{q}-\vec{p}_0)|}{\|\vec{n}\|} = \frac{|n_1 x_Q + n_2 y_Q + n_3 z_Q - d|}{\sqrt{n_1^2 + n_2^2 + n_3^2}}}
        \]
        \item \textbf{Intersección de Dos Planos ($\Pi_1 \cap \Pi_2 = \mathcal{L}$):}
        Como $\vec{v} \perp \vec{n}_1$ y $\vec{v} \perp \vec{n}_2$:
        \[
        \boxed{\vec{v} = \vec{n}_1 \times \vec{n}_2}
        \]
      \end{itemize}
    \end{column}
    \begin{column}{0.42\textwidth}
      \centering
      \begin{tikzpicture}[scale=0.72]
        \fill[navyaccent!15, opacity=0.85] (0,0) -- (3.6,-0.6) -- (5.2,1.5) -- (1.6,2.1) -- cycle;
        \draw[color=navyaccent, thick] (0,0) -- (3.6,-0.6) -- (5.2,1.5) -- (1.6,2.1) -- cycle;
        \node[color=navyaccent, font=\scriptsize\bfseries] at (0.65, 0.25) {$\Pi_1$};

        \fill[orangeacc!15, opacity=0.85] (1.3,-0.9) -- (2.7,2.4) -- (4.1,4.3) -- (2.7,1.0) -- cycle;
        \draw[color=orangeacc, thick] (1.3,-0.9) -- (2.7,2.4) -- (4.1,4.3) -- (2.7,1.0) -- cycle;
        \node[color=orangeacc, font=\scriptsize\bfseries, left=3pt] at (2.2, 3.2) {$\Pi_2$};

        \draw[color=blackmain, line width=2.0pt] (1.6, -0.3) -- (3.5, 2.3);
        \draw[-{Stealth}, color=blackmain, line width=2.2pt] (2.1, 0.38) -- (3.25, 1.95) node[right=4pt, font=\scriptsize\bfseries] {$\vec{v} = \vec{n}_1 \times \vec{n}_2$};

        \draw[-{Stealth}, color=navyaccent, line width=1.3pt] (2.3, 0.65) -- (1.0, 1.8) node[left=2pt, font=\scriptsize\bfseries] {$\vec{n}_1$};
        \draw[-{Stealth}, color=orangeacc, line width=1.3pt] (2.3, 0.65) -- (3.8, 0.35) node[above right=-1pt, font=\scriptsize\bfseries] {$\vec{n}_2$};
      \end{tikzpicture}
    \end{column}
  \end{columns}
\end{frame}

% ---------------------------------------------------------------------
% SLIDE 25: PREGUNTAS DE CIERRE - SÍNTESIS ESTRUCTURAL
% ---------------------------------------------------------------------
\begin{frame}{De la Geometría a la Estructura Lineal}
  \vspace{0.2cm}
  \small
  \begin{enumerate}\itemsep14pt
    \item ¿Cómo se relacionan la descripción interna de un plano (generada por 2 vectores independientes) y la descripción externa (restringida por 1 ecuación ortogonal)? ¿Por qué la suma de parámetros libres y restricciones es la dimensión total del espacio ($2 + 1 = 3$)?
    \item ¿Por qué la intersección de dos planos linealmente independientes en $\R^3$ produce necesariamente una recta? ¿Cómo se generaliza este conteo de codimensión en dimensiones superiores?
    \item ¿En qué sentido la familia de planos paralelos $\vec{n}\cdot\vec{r} = c$ mide desplazamientos de manera dual a un vector flecha ordinario?
  \end{enumerate}
\end{frame}

\end{document}
